\documentclass[a4paper,11pt]{article}
\usepackage[T2A]{fontenc}
\usepackage[utf8]{inputenc}
\usepackage[russian]{babel}
\usepackage{amsmath,amssymb,mathtools}
\renewcommand{\familydefault}{\sfdefault}
\usepackage{geometry}
\geometry{left=23mm,right=23mm,top=23mm,bottom=25mm}
\usepackage{microtype}
\usepackage{booktabs,tabularx,array}
\usepackage{enumitem}
\usepackage{tikz}
\usepackage{hyperref}
\usepackage{xcolor}
\usepackage{fancyhdr}
\usepackage{parskip}
\usepackage{icomma}
\definecolor{accent}{HTML}{087C83}
\definecolor{darkaccent}{HTML}{16484D}
\definecolor{textgray}{HTML}{3B4446}
\definecolor{errorred}{HTML}{A33C39}
\hypersetup{colorlinks=true,linkcolor=accent,pdfauthor={Лёвин Артём Александрович},pdftitle={Ревью домашней работы — Николь Саркисьянц, 08.11.26}}
\setlength{\parindent}{0pt}
\setlength{\parskip}{0.65em}
\setlist[enumerate]{topsep=0.25em,itemsep=0.35em,leftmargin=1.6em}
\pagestyle{fancy}
\fancyhf{}
\fancyhead[L]{\footnotesize\color{textgray}Математическое ожидание и дисперсия}
\fancyhead[R]{\footnotesize\color{textgray}Николь Саркисьянц}
\fancyfoot[C]{\footnotesize\thepage}
\renewcommand{\headrulewidth}{0.3pt}
\setlength{\headheight}{15pt}
\newcommand{\keyidea}{\textcolor{accent}{\bfseries Ключевая идея.}}
\newcommand{\mistake}{\textcolor{errorred}{\bfseries Ошибка.}}
\newcommand{\outcome}{\textcolor{darkaccent}{\bfseries Итог.}}
\newcolumntype{Y}{>{\raggedright\arraybackslash}X}
\begin{document}
\thispagestyle{empty}
\begin{center}
\vspace*{21mm}
{\large\textcolor{darkaccent}{МАТЕМАТИКА}}\par
\vspace{35mm}
{\fontsize{26}{32}\selectfont\bfseries Ревью домашней работы}\par
\vspace{7mm}
{\Large Математическое ожидание, дисперсия\\ и стандартное отклонение}\par
\vspace{25mm}
{\large Николь Саркисьянц}\par
\vspace{8mm}
{\normalsize 8 ноября 2026 г.}\par
\end{center}
\vfill
\begin{center}
\begin{tikzpicture}
\draw[accent!65,line width=0.8pt] (-5,0) .. controls (-2.6,0.65) and (0.3,-0.65) .. (2.3,0) .. controls (3.7,0.4) and (4.5,0.25) .. (5,0.1);
\end{tikzpicture}\par
\vspace{9mm}
{\small Лёвин Артём Александрович\\ эксклюзивно для Николь Саркисьянц}
\end{center}
\clearpage
\section*{Сводная таблица}
В документ включено задание, требующее исправления. Номер в первой колонке ведёт к подробному разбору.\par\vspace{5mm}
\small
\renewcommand{\arraystretch}{1.4}
\begin{tabularx}{\linewidth}{@{}p{11mm}p{24mm}Y p{32mm}p{31mm}@{}}
\toprule
\textbf{№} & \textbf{Тема} & \textbf{Суть ошибки} & \textbf{Ваш ответ} & \textbf{Верный ответ}\\
\midrule
\hyperref[task:10]{\textbf{10}} & Распределение случайной величины & Неверное вычитание при нахождении вероятности $p$; дальнейшие значения вычислены с ошибочным $p$. & $p\approx 0{,}42$; $D=2{,}46$; $\sigma\approx1{,}57$ & $p=0{,}25$; $D=4{,}5$; $\sigma\approx2{,}12$\\
\bottomrule
\end{tabularx}
\normalsize
\clearpage
\thispagestyle{empty}
\vspace*{0.36\textheight}
\begin{center}
{\fontsize{27}{33}\selectfont\bfseries\textcolor{darkaccent}{Разбор ошибок}}
\end{center}
\clearpage
\phantomsection\label{task:10}
\section*{Задание № 10}
\textbf{Условие.} Случайная величина $X$ принимает значения $-2$, $1$ и $4$ с вероятностями $p$, $0{,}5$ и $0{,}5-p$ соответственно. Известно, что $M(X)=1$. Найдите $p$, дисперсию и стандартное отклонение $X$.

\textbf{Ваше решение.} Вы верно составили выражение для математического ожидания:
\[
-2p+0{,}5+4(0{,}5-p)=1.
\]
Далее получено $-6p=-2{,}5$, откуда $p\approx0{,}42$ и $0{,}5-p\approx0{,}08$. С этими вероятностями вычислены $M(X^2)=3{,}46$, $D(X)=2{,}46$ и $\sigma(X)\approx1{,}57$.

\textbf{Ваш ответ.} $p\approx0{,}42$; $D(X)=2{,}46$; $\sigma(X)\approx1{,}57$.

\mistake\ Неверно перенесено постоянное слагаемое при решении линейного уравнения.

\textbf{Где возникает ошибка.} Последний правильный шаг:
\[
-2p+0{,}5+2-4p=1,
\qquad 2{,}5-6p=1.
\]
Первый неверный переход в записи:
\[
\textcolor{errorred}{-6p=-2{,}5}.
\]

\textbf{Почему это неверно.} Чтобы оставить слева только слагаемое с $p$, следует \emph{вычесть $2{,}5$ из обеих частей} уравнения. Правая часть должна стать $1-2{,}5=-1{,}5$. Число $1$ в правой части терять нельзя. Поэтому правильное равенство:
\[
\boxed{-6p=-1{,}5}.
\]

\keyidea\ Сначала аккуратно найдите неизвестную вероятность из условия о математическом ожидании. Затем проверьте, что все вероятности неотрицательны, в сумме дают единицу и воспроизводят заданное ожидание. Для дисперсии используйте уже проверенные вероятности.

\textbf{Исправление с последнего правильного шага.}
\begin{align*}
2{,}5-6p&=1,\\
-6p&=1-2{,}5=-1{,}5,\\
p&=\frac{-1{,}5}{-6}=0{,}25.
\end{align*}
Следовательно, третья вероятность равна $0{,}5-p=0{,}25$.

\textbf{Полное правильное решение.}

\begin{enumerate}
\item \textbf{Проверяем допустимые значения $p$.} Так как $p$ и $0{,}5-p$ являются вероятностями, должны выполняться условия $p\geqslant0$ и $0{,}5-p\geqslant0$. Поэтому $0\leqslant p\leqslant0{,}5$. Полученное $p=0{,}25$ подходит.
\item \textbf{Записываем распределение вероятностей:}
\[
\begin{array}{c|ccc}
X&-2&1&4\\\hline
P&0{,}25&0{,}5&0{,}25
\end{array}
\]
\item \textbf{Находим математическое ожидание квадрата:}
\begin{align*}
M(X^2)&=(-2)^2\cdot0{,}25+1^2\cdot0{,}5+4^2\cdot0{,}25\\
&=4\cdot0{,}25+1\cdot0{,}5+16\cdot0{,}25\\
&=1+0{,}5+4=5{,}5.
\end{align*}
\item \textbf{Находим дисперсию:} по формуле $D(X)=M(X^2)-[M(X)]^2$ и условию $M(X)=1$ получаем
\[
D(X)=5{,}5-1^2=\boxed{4{,}5}.
\]
\item \textbf{Находим стандартное отклонение:}
\[
\sigma(X)=\sqrt{D(X)}=\sqrt{4{,}5}=\frac{3\sqrt2}{2}\approx\boxed{2{,}12}.
\]
\end{enumerate}

\textbf{Проверка.} Сумма вероятностей равна $0{,}25+0{,}5+0{,}25=1$. Кроме того,
\[
M(X)=(-2)\cdot0{,}25+1\cdot0{,}5+4\cdot0{,}25=-0{,}5+0{,}5+1=1.
\]
Условие задачи выполнено. Если подставить записанное Вами приближение $p=0{,}42$, получится $M(X)=2{,}5-6\cdot0{,}42=-0{,}02$ вместо $1$. Это сразу обнаруживает исходную ошибку.

\outcome\ $p=0{,}25$, $D(X)=4{,}5$, $\sigma(X)=\dfrac{3\sqrt2}{2}\approx2{,}12$.

\subsection*{Задача на закрепление}
Случайная величина $Y$ принимает значения $-1$, $2$ и $5$ с вероятностями $q$, $0{,}4$ и $0{,}6-q$ соответственно. Известно, что $M(Y)=2{,}6$. Найдите $q$, дисперсию и стандартное отклонение $Y$.

\vfill
\begin{center}
{\footnotesize\color{textgray}Лёвин Артём Александрович эксклюзивно для Николь Саркисьянц}
\end{center}
\end{document}